\documentclass{article}
\usepackage{TLZLaTeXPack}
\renewcommand{\authorname}{Zhiheng Luo}
\renewcommand{\pdftitle}{Math Olympiad Manifesto}
\renewcommand{\createddate}{May 10, 2026}
\renewcommand{\updateddate}{May 14, 2026}
\renewcommand{\ManifestoRedirctLink}{https://github.com/ZhihengLuo/tlz-latex}

\title{Zhiheng Luo's Math Olympiad Manifesto}
\author{\TLZFullAuthor}

\begin{document}
\maketitle

\newpage
\section{Introduction}
My name is Zhiheng Luo; alternative handles are TLZ or ZhiTuur. The latter is mostly used in gaming contexts.
I am a high school student in Flanders, Belgium, and an aspiring math Olympiad enthusiast.
This document serves as a central repository for my mathematical philosophy and the technical aspects I use in my documents.
The following sections outline the standards I hold for my own work.

\section{Provenance and Navigation}

To maintain transparency in my workflow and help others navigate my resources, this section outlines the entry points to this manifesto and my broader platform.

\subsection{How you got here}
Readers should reach this manifesto via the following ways:
\begin{itemize}
    \item \textbf{GitHub repository:} The primary source for hosting my \texttt{.pdf} and \texttt{.tex} files, as well as the \texttt{TLZLaTeXPack.sty} file used to render this document.
    \item \textbf{Documents and handouts:} Any training material, handouts, problem solves or notes I create will include a hyperlink in the header back to this manifesto.
    \item \textbf{Social links:} Any of my social media profiles in connection with math Olympiad content will contain a link back to this manifesto.
\end{itemize}
The main reason I included this centralised control is if due to saving of my documents you end up with an older version you are always able to find the newest edition.

\subsection{GitHub repository}
As stated before, please consult this \href{https://github.com/ZhiTuur}{\textcolor{red}{link}} for the most up-to-date version of my documents.

\section{Technical}

\subsection{Math notation}
Unless otherwise specified, I adhere to standard IMO notation. 

\subsection{Latex}
\subsubsection{TLZLaTeXPack.sty}
All of my documents are rendered using the \texttt{TLZLaTeXPack.sty} package, which is a custom LaTeX style file I created to ensure consistency across my materials. This package includes settings for fonts, colors, and formatting that align with my personal aesthetic and standards for clarity.
The package is designed to be easily reusable, allowing me to maintain a cohesive look and feel across.
If you are interested in using or adapting the \texttt{TLZLaTeXPack.sty} for your own documents, you can find it in my GitHub repository. Credit is appreciated.
Heavily inspired by Evan Chen's \href{https://github.com/vEnhance/dotfiles/blob/main/texmf/tex/latex/evan/evan.sty}{\textcolor{red}{evan.sty}}.

\subsubsection{Geometry}
I use the standard TikZ library for geometry figures instead of the more programmatic Asymptote. 

\subsection{Pull requests}
If you have suggestions for any of my documents, pull requests are always open on GitHub.
\end{document}
